\documentclass[border=5pt]{standalone}
\usepackage{pgfplots}
\usepackage{pgf-pie} % 饼图：AI 代码可直接使用 \pie
\pgfplotsset{compat=1.18}
\usepgfplotslibrary{fillbetween}  % 面积图 / 堆叠面积图 / \closedcycle 填充路径
% 注：error bars 是 pgfplots 内置功能（在核心 pgfplots.errorbars.code.tex），不需要单独 \usepgfplotslibrary 加载
\usepackage{amsmath}
\usepackage{amssymb}
\usepackage{xcolor}[dvipsnames,svgnames]  % 加载标准色名表（steelblue/teal/orange/coral 等），避免 AI 用预定义色名时报 Undefined color

% 支持中文
\usepackage{fontspec}
\usepackage{xeCJK}
\setCJKmainfont{SimSun}
\setmainfont{Times New Roman}

\begin{document}

\begin{tikzpicture}
\begin{axis}[
    width=13cm,
    height=8cm,
    title={2019--2024年三款产品年销量堆叠面积图},
    xlabel={年份},
    ylabel={销量（百万台）},
    xtick={2019,2020,2021,2022,2023,2024},
    ymin=0,
    ymax=340,
    stack plots=y,
    area style,
    ymajorgrids=true,
    grid style={dashed, gray!30},
    tick align=outside,
    legend style={at={(1.03,0.5)}, anchor=west, draw=black, fill=white}
]
% 底层：产品A（标注为各层原始销量,非累积值）
\addplot[draw=blue!70!black, fill=blue!35, thick, mark=*, mark size=1.5pt,
    nodes near coords, point meta=explicit symbolic,
    every node near coord/.append style={font=\tiny, fill=none, draw=none, inner sep=1pt, anchor=south}]
    coordinates {
        (2019,191.4)[191]
        (2020,206.1)[206]
        (2021,235.3)[235]
        (2022,226.4)[226]
        (2023,231.8)[232]
        (2024,232.0)[232]
    };
% 中层：产品B
\addplot[draw=orange!80!black, fill=orange!40, thick, mark=*, mark size=1.5pt,
    nodes near coords, point meta=explicit symbolic,
    every node near coord/.append style={font=\tiny, fill=none, draw=none, inner sep=1pt, anchor=south}]
    coordinates {
        (2019,49.9)[50]
        (2020,53.2)[53]
        (2021,57.8)[58]
        (2022,60.8)[61]
        (2023,48.5)[49]
        (2024,50.0)[50]
    };
% 顶层：产品C
\addplot[draw=green!50!black, fill=green!35, thick, mark=*, mark size=1.5pt,
    nodes near coords, point meta=explicit symbolic,
    every node near coord/.append style={font=\tiny, fill=none, draw=none, inner sep=1pt, anchor=south}]
    coordinates {
        (2019,18.1)[18]
        (2020,22.6)[23]
        (2021,27.7)[28]
        (2022,28.6)[29]
        (2023,22.7)[23]
        (2024,22.5)[23]
    };
\legend{产品A,产品B,产品C}
\node[anchor=north west, font=\scriptsize] at (axis description cs:0.0,-0.15) {数据来源：IDC/Canalys 2019--2024年全球出货量统计（单位：百万台,标注为各产品当年销量原始值）};
\end{axis}
\end{tikzpicture}

\end{document}